\section{\tr{Operation}}
\subsection{Vario Display}

\tr{The Central Display Information distinguishes between Circling Mode and Straight Flight Mode. If a
turn rate of at least 1°/sec is observed for more than 10 seconds Circling mode is activated. If the
turn rate is less than this value for more than 10 seconds it is switched back to Straight Flight Mode.}

\picnote
	{pictograph-blue-cloud.pdf}
	{\tr{Switching from Variometer Mode to Speed To Fly Mode is independent from switching
	the central display information.}}

\tr{For details see chapter Switching from Circling (Vario) to Speed to Fly.}

\subsubsection{\tr{Circling Mode}}

\includegraphics[width=10cm]{vario-with-explanations.png}

\tr{Content in Circling Mode:}
\begin{itemize}
	\item 	\tr{Info 1: The parameter show is defined in Info 1 Content.}
	\item 	\tr{Data Quality: The colour of the symbol indicates:}
	\begin{itemize}
		\item 	\tr{Green: connection to LARUS Sensor Unit established. Unit has GPS-Fix}
		\item	\tr{Yellow: connection to LARUS Sensor Unit established, Unit has no GPS-Fix}
		\item 	\tr{Red: no connection to LARUS Sensor Unit}
	\end{itemize}
	\item 	\tr{Battery symbol: The colour of the symbol indicates:}
	\begin{itemize}
		\item 	\tr{Green: Voltage of the battery is above Battery Good}
		\item	\tr{Yellow: Voltage of the battery is below Battery Good}
		\item 	\tr{Red: Voltage of the battery is below Battery Low}
	\end{itemize}
	\item 	\tr{Average Climb in Thermal: The average climb rate since starting circling.}
	\item 	\tr{Info 2: The parameter shown is defined in Info 2 Content.}
	\item 	\tr{Current Climb Rate: Time Constant can be changed in Sensor Unit settings Vario TC.}
	\item 	\tr{Average Climb (TC): The average climb rate in the timespan defined in the TC Climb Rate.}
	\item 	\tr{Center Circling Info: Select a Wind or Thermaling Assistance view in Center Circling.}
\end{itemize}


\subsubsection{\tr{Straight Flight Mode}}

\includegraphics[width=10cm]{stf-with-explanations.png}

\tr{Displays that differ from Circling Mode:}
\begin{itemize}
	\item 	\tr{Center Straight Info: Select one of two available graphical representations}
	\item 	\tr{StF Command Band: The position of the band shows whether you are flying too fast or too slowly.
			Positive values mean you are flying too fast and negative values mean you are flying too slow. The
			length of the tape tells you by how much. 1 m/s corresponds to 10 km/h.}
	\item 	\tr{Last Average Climb (TC): The green indicator for the average climb rate is showing the value for the
			last thermal. The value is not updated until the next switch from speed to fly to vario.}
\end{itemize}


\subsubsection{\tr{Available Central Displays}}
\paragraph{\tr{Central Displays in Circling Mode}}

\subparagraph{\tr{Single Arrow}}
\includegraphics[width=5cm]{circling-single-arrow.png}

\tr{The current wind direction is represented by one central arrow. The arrow size is proporional to the wind speed. Changes of the wind direction to a mid term average direction is show by a trailing vane, changes in wind speed to the mid term average speed is represented by the vane’s width. 
The arrow direction is related to North.}


\subparagraph{\tr{Double Arrow}}
\includegraphics[width=5cm]{circling-double-arrow.png}

\tr{Current wind direction and speed is represented by the blue arrow, information related to mid term wind is shown by the grey arrow in the background. Arrow sizes depend on the wind speed.  
The arrow directions are related to North.}


\subparagraph{\tr{Dotted Assistant}}

\includegraphics[width=5cm]{circling-dotted-assistant.png}

\tr{Meaning of the colours of the circle points:}
\begin{itemize}
	\item	\tr{Blue: climb rate is below the average}
	\item	\tr{Red: climb rate is about the average}
	\item	\tr{White/Yellow: climb rate is above the average}
\end{itemize}

The diameter of the dots is proportional to the climb rate.


\subparagraph{\tr{Spider Assistant}}

\includegraphics[width=5cm]{circling-spider-assistant.png}

\tr{Meaning of the colours of the circle points:}
\begin{itemize}
	\item	\tr{Blue: climb rate is below the average}
	\item	\tr{Red: climb rate is about the average}
	\item	\tr{Yellow: climb rate is above the average}
\end{itemize}

\tr{The diameter of the segment is proportional to the climb rate.}



\paragraph{\tr{Central Displays in Straight Flight Mode}}
\subparagraph{\tr{Single Arrow}}

\includegraphics[width=5cm]{straight-single-arrow.png}

\tr{The current wind direction is represented by one central arrow. The arrow size is proporional to the wind speed. Changes of the wind direction to a mid term average direction is show by a trailing vane, changes in wind speed to the mid term average speed is represented by the vane’s width. 
The arrow direction is related to the airplanes longitudinal axis. It is advisable to keep the airplane symbol shown (on) in StF Mode in order to support this details.}


\subparagraph{\tr{Double Arrow}}

\includegraphics[width=5cm]{straight-double-arrow.png}

\tr{Current wind direction and speed is represented by the blue arrow, information related to mid term wind is shown by the grey arrow in the background. Arrow sizes depend on the wind speed.  
The arrow direction is related to the airplanes longitudinal axis. It is advisable to keep the airplane symbol shown in StF Mode in order to support this details.}


\subsubsection{\tr{Switching from Circling to Speed to Fly}}

\tr{Switching the tone information from Speed to Fly to Vario is depending on your Vario Control configuration:}
\begin{itemize}
	\item	\tr{Auto: below 1.1 times speed of best glide ratio: Vario mode, above: Speed To Fly mode. If in Vario mode and still circling, Vario mode keeps active even this speed limit is exceeded.}
	\item	\tr{Input Pin: To use an external switch (at the control stick, flaps lever etc)}
	\item	\tr{SteFly Remote Control Stick}
\end{itemize}


\subsubsection{\tr{Acoustic Quiet Zone in Speed to Fly Mode}}

\tr{You can define a speed range in which the acoustic Speed To Fly signal is quiet.  Set the desired values in StF Upper/Lower Limit in the LARUS Vario Displays advanced settings.}


\subsection{\tr{Control Elements}}

\tr{The device has 2 rotary knobs and a push button that can be used to make various settings.
The following functions are assigned to the button and the two rotary knobs:}
\begin{itemize}
	\item	\tr{Turn the small/upper rotary knob: Volume control}
	\item	\tr{Turn the large/lower rotary knob: Set the MacCready}
	\item	\tr{Short press on the rotary knob: Menu for common settings}
	\item	\tr{Long press on the rotary knob: Menu for configuration settings of the 
			LARUS Vario Display, configuration and calibration of LARUS Sensor Unit}
\end{itemize}

\tr{The following two chapters show the menu structure of the LARUS Vario. Details on the submenu
items are documented in the subsequent chapters..}
